\documentclass[11pt,UTF8,fontset=fandol]{ctexart}

\usepackage[margin=0.82in]{geometry}
\usepackage{amsmath,amssymb}
\usepackage{booktabs,tabularx,array}
\usepackage[dvipsnames]{xcolor}
\usepackage{enumitem}
\usepackage{xurl}
\usepackage{hyperref}

\hypersetup{
  colorlinks=true,
  linkcolor=MidnightBlue,
  citecolor=ForestGreen,
  urlcolor=BrickRed,
  pdftitle={用户证据应该留在上下文，还是写入模型参数？}
}

\newcolumntype{Y}{>{\raggedright\arraybackslash}X}
\newcolumntype{L}[1]{>{\raggedright\arraybackslash}p{#1}}
\renewcommand{\arraystretch}{1.16}
\setlength{\tabcolsep}{5pt}
\setlist[itemize]{leftmargin=1.4em,itemsep=2pt,topsep=3pt}
\setlist[enumerate]{leftmargin=1.6em,itemsep=3pt,topsep=3pt}

\newcommand{\adapter}{\Delta\theta_u}
\newcommand{\memory}{M_u}
\newcommand{\prefscore}{U_{\mathrm{pref}}}
\newcommand{\taskscore}{U_{\mathrm{task}}}

\begin{document}

\begin{center}
  {\LARGE\bfseries 用户证据应该留在上下文，还是写入模型参数？\par}
  \vspace{5pt}
  {\large 个性化模型的方法设计笔记\par}
  \vspace{4pt}
  {\small 2026 年 8 月 29 日}
\end{center}

\section{结论}

新出现的用户偏好，先放在提示词或外部记忆里。这样容易查看、修改和删除。

只有满足下面几项时，才值得训练用户适配器（adapter）：
\begin{itemize}
  \item 同一偏好在多个会话中反复出现；
  \item 偏好长期稳定；
  \item 偏好会影响多类任务；
  \item 即使不提供用户档案，偏好也应该生效。
\end{itemize}

参数更新不是默认答案。它必须比上下文更有效，同时不能降低任务成功率。

本文提出一个简单实验。对同一批用户证据，做两个版本：
\begin{enumerate}
  \item \textbf{上下文版}：把证据保存在提示词或外部记忆中，不改模型参数；
  \item \textbf{参数版}：用同一批证据训练用户适配器，测试时不再提供这批证据。
\end{enumerate}
两个版本使用同一个起点，并在同一组后续任务上比较。只有参数版的个性化效果更好，且任务成功率不低于上下文版，才选择参数版。否则保留上下文。

这仍是一个待验证的方案，不是已经成立的新算法。

\section{要区分的三种信息}

\begin{table}[ht]
  \centering
  \caption{三种信息的作用不同。}
  \begin{tabularx}{\textwidth}{@{}L{0.18\textwidth}L{0.35\textwidth}Y@{}}
    \toprule
    信息 & 用法 & 特点 \\
    \midrule
    提示词或外部记忆 &
    在运行时提供用户事实、偏好或历史记录。 &
    容易修改，但会占用上下文，也可能检索错误。 \\
    训练提示 &
    训练时给教师模型的提示，用来生成训练目标。 &
    它是训练信号，不一定是部署时的用户状态。 \\
    用户适配器 \(\adapter\) &
    针对用户 \(u\) 学到的一组参数。 &
    不需要重复输入用户档案，但更难检查和修改。 \\
    \bottomrule
  \end{tabularx}
\end{table}

这里的在线更新指部署后更新。系统先服务用户，再用新交互更新参数，并得到新的适配器。用固定数据集训练、只是在训练时重新采样，并不等于部署后学习。

本文讨论的个性化智能体还必须完成任务。只改变文字风格，不等于个性化智能体。

\section{这些设计从哪里来}

当前代码可能有错误，也可能继续变化，因此不作为方法依据。本文只使用论文中已经说明的设计和结果。

\begin{table}[ht]
  \centering
  \caption{设计依据。}
  \begin{tabularx}{\textwidth}{@{}L{0.20\textwidth}L{0.52\textwidth}Y@{}}
    \toprule
    设计 & 依据 & 本文如何使用 \\
    \midrule
    按时间划分数据 &
    SPRInG 把每个用户的记录按时间分成五个阶段。每个阶段前 \(90\%\) 用于更新，后 \(10\%\) 用于评估 \cite{kim2026spring}。 &
    直接采用 \\
    比较上下文和参数 &
    OPPU 比较检索、用户档案提示和每用户 PEFT
    \cite{tan2024oppu}。SPRInG 也分别评估适配器和检索。 &
    作为必要对照 \\
    比较外部经验和参数学习 &
    EvolveR 的主要方法把经验当作上下文。它还测试了把经验写入参数的版本
    \cite{wu2025evolver}。 &
    作为研究动机 \\
    任务成功率 &
    EvolveR 用最终答案是否正确作为任务结果奖励。我们的智能体也必须先完成任务。 &
    作为独立指标 \\
    安全检查 &
    SPRInG 只把安全过滤和投毒防护列为未来工作。 &
    不放进核心方法 \\
    置信区间 &
    SPRInG 用它报告实验结果，但不用它决定每次更新。 &
    只用于结果报告 \\
    回退 &
    上述论文没有给出可直接继承的回退方法。 &
    暂不加入 \\
    \bottomrule
  \end{tabularx}
\end{table}

\section{现有方法已经做了什么}

\paragraph{OPPU。}
OPPU 用每个用户的历史数据训练一个 PEFT 模块，也可以同时使用检索或用户档案
\cite{tan2024oppu}。它是离线方法，不处理部署后的逐次更新。

\paragraph{Personalized DPO。}
Personalized RLHF 从离线偏好数据学习用户表示和 LoRA 参数
\cite{li2024prlhf}。它没有研究外部记忆，也没有研究智能体的任务结果。

\paragraph{SPRInG。}
SPRInG 是最接近的工作 \cite{kim2026spring}。它同时使用用户 LoRA 和检索
缓冲区，并按时间持续更新。因此，``外部记忆加适配器'' 不是新贡献。

SPRInG 根据似然变化选择训练样本。更新后仍然难学的样本进入缓冲区。它没有对同一批证据分别训练上下文版和参数版，再用后续任务选择其中一个。

\paragraph{EvolveR。}
EvolveR 让智能体从轨迹中总结经验，并把经验放在外部经验库（Experience Base）中
\cite{wu2025evolver}。它还测试了直接学习这些经验。论文报告的平均分从
\(0.382\) 降到 \(0.371\)。这说明，未经筛选就把经验写入参数，可能带来噪声。

\paragraph{CoSteer。}
CoSteer 不训练用户参数 \cite{lv2025costeer}。本地小模型分别读取和不读取用户
上下文，两次输出的差值用来调整云端模型。它说明运行时方法也可能有效，因此应作为对照。

\section{我们要验证的方法}

设用户 \(u\) 当前有适配器 \(\adapter^{(t)}\) 和外部记忆 \(\memory^{(t)}\)。
新证据 \(E_t\) 来自该用户之后的交互。

\paragraph{上下文版。}
\begin{equation}
  M_{u,\mathrm{ctx}}
  =\operatorname{MemoryUpdate}\!\left(\memory^{(t)},E_t\right),
  \qquad
  \Delta\theta_{u,\mathrm{ctx}}=\adapter^{(t)}.
\end{equation}

\paragraph{参数版。}
\begin{equation}
  \Delta\theta_{u,\mathrm{par}}
  =\operatorname{TrainAdapter}\!\left(\adapter^{(t)},E_t\right),
  \qquad
  M_{u,\mathrm{par}}=\memory^{(t)}.
\end{equation}

参数版和上下文版必须使用同一批证据、同一个起点和同一组后续任务
\(V_t\)。参数版测试时不能再把 \(E_t\) 放进提示词，否则无法知道参数学到了什么。

选择规则只有两项：
\begin{equation}
  \begin{aligned}
  \prefscore(m_{\mathrm{par}};V_t)
    &> \prefscore(m_{\mathrm{ctx}};V_t),\\
  \taskscore(m_{\mathrm{par}};V_t)
    &\ge \taskscore(m_{\mathrm{ctx}};V_t).
  \end{aligned}
\end{equation}

两项都满足，选择参数版。否则选择上下文版。若证据互相矛盾或数量太少，暂不更新。

这条规则是本文提出的实验方案，不是从现有论文直接复制的结论。

\section{实验怎么做}

\subsection{按时间划分}

先按完整会话排序，再提取用户档案或生成训练数据：
\begin{equation}
  H_u^{\mathrm{past}}
  \prec E_t
  \prec V_t
  \prec H_u^{\mathrm{test}}.
\end{equation}
过去数据用于初始化，\(E_t\) 用于生成两个版本，\(V_t\) 用于选择，最后一段只用于测试。

\subsection{最少需要的对照}

\begin{table}[ht]
  \centering
  \caption{实验对照。}
  \begin{tabularx}{0.94\textwidth}{@{}L{0.32\textwidth}Y@{}}
    \toprule
    版本 & 回答的问题 \\
    \midrule
    基础模型 & 不做个性化时表现如何？ \\
    基础模型 + 提示词/记忆 & 只用上下文能做到什么？ \\
    仅适配器 & 参数是否真的学到用户偏好？ \\
    适配器 + 提示词/记忆 & 两种方式能否互补？ \\
    OPPU & 普通每用户 PEFT 是否足够？ \\
    SPRInG & 已有的持续更新加检索是否足够？ \\
    CoSteer（条件允许时） & 不训练参数能否达到相近效果？ \\
    本文方案 & 按证据选择保存位置是否更好？ \\
    \bottomrule
  \end{tabularx}
\end{table}

至少分别报告三类结果：
\begin{itemize}
  \item 后续任务中的个性化效果；
  \item 环境可以验证的任务成功率；
  \item 提示词、训练和推理成本。
\end{itemize}
不要把它们合成一个总分。置信区间只用于说明结果是否稳定，不参与每次更新。

\section{创新边界}

以下内容已有工作做过，不能作为新贡献：
\begin{itemize}
  \item 每个用户一个适配器；
  \item 持续更新用户适配器；
  \item 同时使用外部记忆和适配器；
  \item 比较外部上下文和参数学习。
\end{itemize}

可能的新点更窄：对同一批证据，生成上下文版和参数版，再用同一用户的后续任务决定证据应该放在哪里。

这个差异仍需更完整的文献检查和实验验证。当前只能称为研究假设。

\section{下一步}

\begin{enumerate}
  \item 选择带时间信息的同一用户数据。
  \item 根据论文或官方代码实现 OPPU、SPRInG 和提示词/记忆对照。
  \item 固定一种适配器训练方法和一种外部记忆写入方法。
  \item 先跑基础模型、仅上下文、仅适配器和两者结合。
  \item 再跑本文的配对比较。
  \item 个性化和任务成功都通过后，再扩展到真实工具调用的智能体任务。
\end{enumerate}

\begin{thebibliography}{9}

\bibitem{tan2024oppu}
Zhaoxuan Tan, Qingkai Zeng, Yijun Tian, Zheyuan Liu, Bing Yin, and Meng Jiang.
\newblock Democratizing Large Language Models via Personalized
Parameter-Efficient Fine-Tuning.
\newblock EMNLP, 2024.
\newblock \url{https://aclanthology.org/2024.emnlp-main.372/}.

\bibitem{li2024prlhf}
Xinyu Li, Ruiyang Zhou, Zachary C. Lipton, and Leqi Liu.
\newblock Personalized Language Modeling from Personalized Human Feedback.
\newblock arXiv:2402.05133, 2024.
\newblock \url{https://arxiv.org/abs/2402.05133}.

\bibitem{kim2026spring}
Seoyeon Kim and Jaehyung Kim.
\newblock SPRInG: Continual LLM Personalization via Selective Parametric
Adaptation and Retrieval-Interpolated Generation.
\newblock arXiv:2601.09974, 2026.
\newblock \url{https://arxiv.org/abs/2601.09974}.

\bibitem{wu2025evolver}
Rong Wu, Xiaoman Wang, Jianbiao Mei, and others.
\newblock EvolveR: Self-Evolving LLM Agents through an Experience-Driven
Lifecycle.
\newblock arXiv:2510.16079, 2025.
\newblock \url{https://arxiv.org/abs/2510.16079}.

\bibitem{lv2025costeer}
Hang Lv, Sheng Liang, Hao Wang, and others.
\newblock CoSteer: Collaborative Decoding-Time Personalization via Local
Delta Steering.
\newblock arXiv:2507.04756, 2025.
\newblock \url{https://arxiv.org/abs/2507.04756}.

\end{thebibliography}

\end{document}
